% ================================================================
% ==== FILE: content/k_levels/k7.tex
% ================================================================

% ==============================
%  Ontology of Continua — Core
%  K-level Module: K7 (Social Systems)
%  FULL MODULE — FINAL
% ==============================



\paragraph{Canonical tuple projection note.}
Local K-level displays in this module are projections of the canonical public tuple \(K=(\Omega,\partial\Omega,A,\Theta,P,J,C,k,M)\). When a display omits \(M\), reorders components, or uses level-indexed shorthand, the omitted embedding context is inherited from the surrounding level discussion rather than introduced as a competing definition.

\subsubsection{\texorpdfstring{$K_7$}{K_7} Übersicht}
\label{sec:k7-overview}

$K_7$ ist die Ebene der \emph{sozialen Kontinua}: Systeme, deren Zustände,
Achsen und Potentiale entstehen durch \emph{Wechselwirkung zwischen mehreren $K_6$
kognitive Kontinua}. Es ist die kleinste Ebene, auf der:

\begin{itemize}
    \item gemeinsame Modelle und gemeinsames Wissen entstehen,
    \item Kommunikation bildet eine strukturelle Dimension,
    \item Kooperation und Konflikt werden zu dynamischen Kräften,
    \item soziale Identität und Rollenbildung treten auf,
    \item kollektives Gedächtnis und Normen stabilisieren das Verhalten,
    \item Zyklen auf Gruppenebene definieren die zeitliche Organisation.
\end{itemize}

$K_7$ ist keine Summe von Agenten: Es ist ein neues Kontinuum mit seinem eigenen
zulässige Zustände, Schwellenwerte, Flüsse und Todesmechanismen.

% ================================================================
\subsubsection{Zustandsraum \texorpdfstring{$\Omega(K_7)$}{\Omega(K_7)}}

Ein Zustand des sozialen Kontinuums umfasst:

\[
\Omega(K_7)=
\{
G_{\mathrm{soc}},\;
M_{\mathrm{shared}},\;
C_{\mathrm{comm}},\;
R_{\mathrm{Rollen}},\;
N_{\mathrm{Normen}},\;
H_{\mathrm{hist}},\;
S_{\mathrm{struc}},\;
\mu^{(7)}_t,\;
\Sigma_{\mathrm{Identität}}
\}.
\]

Komponenten:
\begin{itemize}
    \item $G_{\mathrm{soc}}$ – sozialer Graph (Agenten als Knoten, Bindungen als Kanten),
    \item $M_{\mathrm{shared}}$ – gemeinsame Modelle (kollektive Überzeugungen),
    \item $C_{\mathrm{comm}}$ — Kommunikationsstruktur,
    \item $R_{\mathrm{roles}}$ — Rollenzuweisung (Status, Funktionen),
    \item $N_{\mathrm{norms}}$ – Normen und Einschränkungen,
    \item $H_{\mathrm{hist}}$ – kollektive Geschichten und Erzählungen,
    \item $S_{\mathrm{struc}}$ — Strukturpositionen und Module,
    \item $\mu^{(7)}_t$ — Verteilung der Gruppenzustände über die Zeit,
    \item $\Sigma_{\mathrm{identity}}$ – soziale Identitätsstrukturen.
\end{itemize}

Die Zulässigkeit wird durch $M_7$ und Gruppenkohärenzschwellen eingeschränkt.

% ================================================================
\subsubsection{Grenze \texorpdfstring{$\partial\Omega(K_7)$}{\partial\Omega(K_7)}}

Die Grenze umfasst:
\begin{itemize}
    \item Zusammenbruch gemeinsam genutzter Modelle,
    \item Fragmentierung des sozialen Graphen,
    \item Aufschlüsselung der Kommunikation,
    \item Verlust der normativen Kohärenz,
    \item-Rolleninstabilität,
    \item Identitätszusammenbruch,
    \item Unfähigkeit, das kollektive Gedächtnis zu stabilisieren.
\end{itemize}

Das Überschreiten von $\partial\Omega(K_7)$ zerstört die Gruppenkontinuität.

% ================================================================
\subsubsection{Achsen \texorpdfstring{$A(K_7)$}{A(K_7)}}

Aus dem Kernmodul „Soziale Systeme“:

\[
A(K_7)=
\{
A_{\mathrm{comm}},\;
A_{\mathrm{coop}},\;
A_{\mathrm{conf}},\;
A_{\mathrm{norm}},\;
A_{\mathrm{Rolle}},\;
A_{\mathrm{id}},\;
A_{\mathrm{gemeinsames Modell}}
\}.
\]

Interpretation:
\begin{itemize}
    \item $A_{\mathrm{comm}}$ — Kommunikationsachse (Kanäle, Codes),
    \item $A_{\mathrm{coop}}$ – Kooperationsachse,
    \item $A_{\mathrm{conf}}$ – Konflikt-/Konkurrenzachse,
    \item $A_{\mathrm{norm}}$ – Normativitätsachse,
    \item $A_{\mathrm{role}}$ — Rollen-Positions-Achse,
    \item $A_{\mathrm{id}}$ – soziale Identitäten,
    \item $A_{\mathrm{shared-model}}$ – kollektive Modelle/Überzeugungen.
\end{itemize}

Diese Achsen können nicht auf $K_6$ kognitive Unterschiede reduziert werden.

% ================================================================
\subsubsection{Potentiale \texorpdfstring{$P(K_7)$}{P(K_7)}}

Potenziale treiben gesellschaftliche Dynamiken:

\[
P(K_7)=
\{
P_{\mathrm{comm}},\;
P_{\mathrm{coop}},\;
P_{\mathrm{conf}},\;
P_{\mathrm{norm}},\;
P_{\mathrm{id}},\;
P_{\mathrm{Rolle}},\;
P_{\mathrm{shared}}
\}.
\]

Bedeutung:
\begin{itemize}
    \item $P_{\mathrm{comm}}$: Kommunikationseffektivität,
    \item $P_{\mathrm{coop}}$: Kooperationsvorteil,
    \item $P_{\mathrm{conf}}$: Konfliktdruck,
    \item $P_{\mathrm{norm}}$: normativer Zusammenhalt,
    \item $P_{\mathrm{id}}$: Identitätsstabilisierungspotential,
    \item $P_{\mathrm{role}}$: funktionaler Rollendruck,
    \item $P_{\mathrm{shared}}$: Kohärenzpotential des gemeinsamen Modells.
\end{itemize}

% ================================================================
\subsubsection{Schwellenwerte \texorpdfstring{$\Theta(K_7)$}{\Theta(K_7)}}

\[
\Theta(K_7)=
\{
\Theta_{\mathrm{comm}},\;
\Theta_{\mathrm{coh}},\;
\Theta_{\mathrm{norm}},\;
\Theta_{\mathrm{id}},\;
\Theta_{\mathrm{Rolle}},\;
\Theta_{\mathrm{shared}},\;
\Theta_{\mathrm{Fragment}}
\}.
\]

Interpretation:
\begin{itemize}
    \item $\Theta_{\mathrm{comm}}$: minimale Kommunikationsbandbreite,
    \item $\Theta_{\mathrm{coh}}$: Kohärenzschwelle für Gruppenidentität,
    \item $\Theta_{\mathrm{norm}}$: minimale normative Ausrichtung,
    \item $\Theta_{\mathrm{id}}$: Identitäts-Persistenz-Schwelle,
    \item $\Theta_{\mathrm{role}}$: Rollenstabilitätsschwelle,
    \item $\Theta_{\mathrm{shared}}$: Kohärenzschwelle für gemeinsame Modelle,
    \item $\Theta_{\mathrm{fragment}}$: Fragmentierungsschwelle des sozialen Diagramms.
\end{itemize}

Das Überschreiten eines dieser Punkte kann zum Zusammenbruch von $K_7$ führen.

% ================================================================
\subsubsection{Flüsse \texorpdfstring{$J(K_7)$}{J(K_7)}}

\begin{itemize}
    \item $J_{\mathrm{comm}}$ — Kommunikationsflüsse,
    \item $J_{\mathrm{coop}}$ — Kooperationsdynamik,
    \item $J_{\mathrm{conf}}$ — Konflikteskalation oder Deeskalation,
    \item $J_{\mathrm{norm}}$ – normaktualisierende Flüsse,
    \item $J_{\mathrm{id}}$ — Identitätsübergänge,
    \item $J_{\mathrm{role}}$ – Rollenneuzuweisung,
    \item $J_{\mathrm{shared}}$ – Flussaktualisierung gemeinsam genutzter Modelle,
    \item $J_{\mathrm{PE}\to\mathrm{comm}}$: Projektion von Vorhersagefehlersignalen von $K_6$ in die Kommunikation,
    \item $J_{\mathrm{comm}\to A_{\mathrm{error}}}$: Feedback aus der sozialen Kommunikation zurück in kognitive Vorhersageachsen.
\end{itemize}

Stabilität erfordert:
\[
\sigma(J(K_7)) < \sigma_{\max}(K_7).
\]

% ================================================================
\subsubsection{Zyklen \texorpdfstring{$C(K_7)$}{C(K_7)}}

\[
C(K_7)=
\{
C_{\mathrm{comm}},\;
C_{\mathrm{coop}},\;
C_{\mathrm{conf}},\;
C_{\mathrm{norm}},\;
C_{\mathrm{id}},\;
C_{\mathrm{shared}}
\}.
\]

Beschreibungen:
\begin{itemize}
    \item $C_{\mathrm{comm}}$: Kommunikation → Interpretation → Antwort → Aktualisierung,
    \item $C_{\mathrm{coop}}$: Kooperationszyklen (Vertrauens-Verstärkungs-Schleifen),
    \item $C_{\mathrm{conf}}$: Konfliktlösungszyklen,
    \item $C_{\mathrm{norm}}$: Normerstellung, Durchsetzung, Überarbeitung,
    \item $C_{\mathrm{id}}$: Identitätsbildung, Aushandlung, Stabilisierung,
    \item $C_{\mathrm{shared}}$: Aufrechterhaltung gemeinsamer Überzeugungen und des kollektiven Gedächtnisses.
\end{itemize}

Jeder Zyklus erzeugt eine soziale Zeitstruktur.

% ================================================================
\subsubsection{Zeit \texorpdfstring{$\tau(K_7)$}{\tau(K_7)}}

Soziale Zeit entsteht aus stabilen Zyklen:

\[
\tau(K_7)
= \min_j \tau_{C_j},
\qquad
\Pi(C_j) > \Theta_{\mathrm{Zeit}}.
\]

Die Zeit von $K_7$ ist typischerweise langsamer und träger als die Zeit von $K_6$
kollektive Trägheit.

% ================================================================
\subsubsection{Kontinuität \texorpdfstring{$k(K_7)$}{k(K_7)}}

Aus dem Modul „Soziale Systeme“:

\[
k_7 =
H(\Omega(K_7))
\cdot
\frac{|C_{\max}^{(7)}|}{|C_{\mathrm{all}}|}
\cdot
\frac{|A_{\mathrm{aktiv}}|}{|A_{\max}|}
\cdot
\left(1-\frac{T_7}{\Theta_7}\right)_+
\cdot
\left(1-\frac{\sigma(J)}{\sigma_{\max}}\right).
\]

Wo:
\begin{itemize}
    \item $C_{\max}^{(7)}$ — größte kohärente soziale Komponente,
    \item $T_7$ – strukturelle Spannung des sozialen Kontinuums,
    \item $\Theta_7$ – dominanter Schwellenwert (Kohärenz- oder Shared-Model-Grenze),
    \item $\sigma(J)$ – Volatilität sozialer Ströme.
\end{itemize}

Hohe Fragmentierung $\Rightarrow k_7\to 0$.

% ================================================================
\subsubsection{Strukturelle Spannung \texorpdfstring{$T(K_7)$}{T(K_7)}}

Quellen:
\begin{itemize}
    \item Kommunikationsfehler,
    \item-Normkonflikte,
    \item-Rolleninstabilität,
    \item Fragmentierung gemeinsam genutzter Modelle,
    \item Identitätskonflikte,
    \item schnelle Verschiebungen im Kooperations-/Konfliktgleichgewicht,
    \item Überkonzentration des Einflusses oder strukturelle Engpässe.
\end{itemize}

Zusammenbrechen, wenn:
\[
T(K_7) > \Theta_{\mathrm{coh}}.
\]

% ================================================================
\subsubsection{Energie \texorpdfstring{$E(K_7)$}{E(K_7)}}

\[
E(K_7)
=
E_{\mathrm{comm}}
+ E_{\mathrm{coop}}
+ E_{\mathrm{conf}}
+ E_{\mathrm{norm}}
+ E_{\mathrm{id}}
+ E_{\mathrm{shared}}
+ E_{K_6\bis K_7}.
\]

Interpretation:
\begin{itemize}
    \item Kommunikationskosten,
    \item Kosten für die Aufrechterhaltung der Zusammenarbeit,
    \item Energie der Konflikteskalation,
    \item normative Durchsetzungskosten,
    \item Identitätspflege,
    \item Aufrechterhaltung gemeinsamer Überzeugungen,
    \item Kopplungsenergie, geerbt von $K_6$.
\end{itemize}

% ================================================================
\subsubsection{Operatoren auf \texorpdfstring{$K_7$}{K_7} (\texorpdfstring{$\Psi$}{\Psi}, \texorpdfstring{$\Phi$}{\Phi}, \texorpdfstring{$\Lambda$}{\Lambda}, \texorpdfstring{$U$}{U}, \texorpdfstring{$\Chi$}{\Chi})}

\begin{itemize}
    \item $\Psi_{7\to8}$: Geburt institutioneller und zivilisatorischer Kontinua,
    \item $\Phi$: Entwicklung sozialer Strukturen und Rollen,
    \item $\Lambda$: Vereinigung von Subsystemen zu zusammenhängenden Gruppen,
    \item $U$: Stabilisierung von Normen und Identitäten,
    \item $\Chi$: Verzweigung in Subkulturen, Fraktionen und spezialisierte Gruppen.
\end{itemize}

% ================================================================
\subsubsection{Prozesse auf \texorpdfstring{$K_7$}{K_7}}

Typisch:
\begin{itemize}
    \item Kommunikation,
    \item Kooperation und Koalitionsbildung,
    \item Konflikteskalation und -lösung,
    \item Normbildung, Durchsetzung und Überarbeitung,
    \item Rollenverteilung und Autoritätsbildung,
    \item Identitätserstellung und -transformation,
    \item Bildung gemeinsamer Erzählungen und kollektiven Gedächtnisses,
    \item soziales Lernen und Verbreitung.
\end{itemize}

% ================================================================
\subsubsection{Vorhersagen für \texorpdfstring{$K_7$}{K_7}}

\begin{enumerate}
    \item Soziale Kontinua brechen zusammen, wenn die Kohärenz des gemeinsamen Modells unter $\Theta_{\mathrm{shared}}$ fällt.
    \item Kommunikationsengpässe lassen auf eine Fragmentierung schließen.
    \item Normative Instabilität geht dem Zusammenbruch der Identität voraus.
    \item Hochmodulare Gruppen widerstehen Lärm und Konflikten.
    \item Der Übergang zu $K_8$ erfordert anhaltende normative und kommunikative Stabilität.
\end{enumerate}

% ================================================================
\subsubsection{Experimente für \texorpdfstring{$K_7$}{K_7}}

Mögliche empirische Analoga:
\begin{itemize}
    \item Multiagenten-Lernsimulationen zur Verstärkung,
    \item Netzwerkfragmentierungsstudien,
    \item Normdiffusionsexperimente,
    \item Tests zur Stabilität des kollektiven Gedächtnisses,
    \item Experimente zur Störung der Kommunikationsbandbreite,
    \item Identitätskohärenzmessungen in sozialen Gruppen.
\end{itemize}

% ================================================================
\subsubsection{Zusammenbruch und Tod von \texorpdfstring{$K_7$}{K_7}}

Der Tod tritt ein, wenn:
\[
\Omega(K_7) = \varnothing.
\]

Mechanismen:
\begin{itemize}
    \item Fragmentierung des sozialen Diagramms,
    \item Verlust der Kommunikationsfähigkeit,
    \item normative Desintegration,
    \item Identitätszusammenbruch,
    \item Wettbewerbsdynamik,
    \item Unfähigkeit, gemeinsam genutzte Modelle zu stabilisieren.
\end{itemize}

% ================================================================
\subsubsection{Falschbarkeit von \texorpdfstring{$K_7$}{K_7}}

Die Theorie sagt voraus:
\begin{itemize}
    \item Existenz von Kohärenzschwellen,
    \item messbare Auswirkungen auf die Kommunikationsbandbreite,
    \item klarer Zusammenhang zwischen Fragmentierung und Zusammenbruch,
    \item stabile soziale Zyklen mit charakteristischen Perioden,
    \item Notwendigkeit von Shared-Model-Strukturen für Gruppenpersistenz.
\end{itemize}

Eine Fälschung würde sich aus der Beobachtung einer großräumigen sozialen Stabilität ergeben
ohne diese Strukturen.

% ================================================================
\subsubsection{Verzweigung / ontologische Position von \texorpdfstring{$K_7$}{K_7}}

Filialen:
\begin{itemize}
    \item Genossenschaften vs. Konkurrenzgesellschaften,
    \item stark normative vs. schwach normative Gruppen,
    \item identitätszentrierte vs. rollenzentrierte Strukturen,
    \item zentralisierte vs. dezentrale Kommunikationsnetzwerke.
\end{itemize}

Ontologischer Ort:
\[
K_6 \to K_7 \to K_8.
\]

% ================================================================
\subsubsection{Beziehung zu M-Räumen}

$M_7$ Einschränkungen:
\begin{itemize}
    \item Social-Graph-Topologien,
    \item Bereich zulässiger Normsysteme,
    \item Kommunikationskomplexität,
    \item Identitätskohärenz,
    \item strukturelle Modularität,
    \item-Kopplung zu $M_6$ (kognitive Ebene).
\end{itemize}

Die Kompatibilität mit $M_7$ bestimmt, ob ein soziales Kontinuum existieren kann.
